\originalchapter{期望的线性性、次序统计量与协方差}
\sourcelecture{23--24}
计算一个和的平均值, 通常不必先求这个和的完整分布. 关键是把复杂计数拆成容易求平均值的小量. 但平均值相加与方差相加有不同条件: 前者不要求独立, 后者必须检查交叉项. 本章先建立这一区别, 再讨论排序后的样本及相关性.

\originalsection{期望的线性性与尾概率}
\begin{theorem}[线性性]
若 $X_1,\ldots,X_n$ 可积, 即 $\E|X_i|<\infty$, 而 $a_1,\ldots,a_n$ 是实常数, 则
\[\E\left[\sum_{i=1}^n a_iX_i\right]=\sum_{i=1}^n a_i\E X_i.\]
这里允许各随机变量彼此依赖.
\end{theorem}
\begin{proof}
离散情形中, 设两变量联合质量为 $p(x,y)$. 因为
$\sum_{x,y}|ax+by|p(x,y)\le |a|\E|X|+|b|\E|Y|<\infty$, 可交换绝对收敛级数的求和顺序. 因而
\[
\begin{aligned}
\E(aX+bY)&=\sum_{x,y}(ax+by)p(x,y)\\
&=a\sum_x x\sum_y p(x,y)+b\sum_y y\sum_x p(x,y)\\
&=a\E X+b\E Y.
\end{aligned}
\]
若有联合密度, 同样把求和换成二重积分; 绝对可积保证交换次序合法. 对任意类型的变量, 期望是对同一概率空间的积分, 积分的有限线性性给出相同结论. 重复二变量结论即得有限和公式. 这里使用的是积分的线性性, 而不是联合分布的分解.
\end{proof}

\begin{definition}[指标变量]
事件 $A$ 的指标变量 $\ind_A$ 在 $A$ 发生时取1, 否则取0. 它的期望为 $\E\ind_A=\Prob(A)$.
\end{definition}
若 $N$ 统计一组事件发生的次数, 写成 $N=\sum_i\ind_{A_i}$ 后便得到 $\E N=\sum_i\Prob(A_i)$. 事件间是否独立, 对这一结论没有影响.
\begin{example}
把 $m$ 个编号包裹随机一一分配给 $m$ 个地址, 所有排列等可能. 求送对地址的包裹数 $N$ 的期望.
\end{example}
\begin{solution}
令 $I_i$ 表示第 $i$ 个包裹送到第 $i$ 个地址. 每个包裹的目的地在 $m$ 个地址中均匀分布, 故 $\E I_i=1/m$. 于是 $\E N=m(1/m)=1$. 这个平均数不随规模变化; 它并不表示每次恰好送对一个.
\end{solution}

\begin{theorem}[尾积分公式]
若 $X\ge0$, 则允许两边同为无穷的恒等式
\[\E X=\int_0^\infty\Prob(X>t)\dd t\]
成立. 若 $X$ 是非负整数变量, 则 $\E X=\sum_{k=0}^\infty\Prob(X>k)$.
\end{theorem}
\begin{proof}
每个固定 $x\ge0$ 都满足 $x=\int_0^\infty\ind_{\{t<x\}}\dd t$. 把 $x$ 换成 $X$, 非负积分可交换顺序, 得
$\E X=\int_0^\infty\E\ind_{\{X>t\}}\dd t$. 非负交换可由先限制到有限区间、再用单调逼近理解, 无须假设期望有限. 整数情形逐点有 $X=\sum_{k\ge0}\ind_{\{X>k\}}$, 用同样的非负交换得到级数式.
\end{proof}
几何上, 平均值是生存函数 $1-F_X(t)$ 下的面积. 这也解释逆分布抽样: 若 $U$ 在 $(0,1)$ 均匀分布, 定义 $Q(u)=\inf\{x:F_X(x)\ge u\}$, 则 $\Prob(Q(U)\le x)=F_X(x)$. 分布函数有跳跃或平坦部分时应使用这个广义逆, 不能假定普通反函数总存在.
\begin{example}
若等待时间 $T$ 满足 $\Prob(T>t)=\ee^{-\lambda t}$, $t\ge0$, 求其平均值.
\end{example}
\begin{solution}
直接积分得 $\E T=\int_0^\infty\ee^{-\lambda t}\dd t=1/\lambda$. 此处不需先求密度, 尾概率本身已经足够.
\end{solution}

\originalsection{排序后的样本}
\begin{definition}[次序统计量]
将独立同分布的 $X_1,\ldots,X_n$ 从小到大排列为 $X_{(1)}\le\cdots\le X_{(n)}$, 称这些变量为次序统计量. 若共同分布有密度, 两个样本相等的概率为0, 故排序几乎必然严格.
\end{definition}
设共同分布函数为 $F$, 密度为 $f$. 最大值不超过 $x$ 等价于所有样本不超过 $x$, 因而
\[\Prob(X_{(n)}\le x)=F(x)^n,\qquad
\Prob(X_{(1)}>x)=(1-F(x))^n.\]
对这些函数求导, 得最大值密度 $nf(x)F(x)^{n-1}$, 最小值密度 $nf(x)(1-F(x))^{n-1}$. 若样本均匀分布于 $(0,1)$, 最大值密度为 $nx^{n-1}$, 平均值为 $n/(n+1)$.

\begin{theorem}[联合排序密度与随机名次]
在 $x_1<\cdots<x_n$ 上, 排序向量的联合密度是
\[n!\prod_{i=1}^nf(x_i),\]
区域外为0. 记录各原始编号对应名次的排列, 在 $n!$ 个排列中均匀分布, 且与排序向量独立.
\end{theorem}
\begin{proof}
严格有序的一个小区域有 $n!$ 个原像, 分别对应原始坐标的各排列. 每个原像的密度均为 $\prod_i f(x_i)$, 坐标置换的体积因子为1, 所以总密度乘 $n!$. 更精确地, 对有序区域中的任一集合 $B$ 和任一固定排列 $\pi$, 联合概率为 $\int_B\prod_i f(x_i)\dd x_1\cdots\dd x_n$, 而排序向量落在 $B$ 的概率为其 $n!$ 倍. 所以联合概率等于 $(1/n!)\Prob((X_{(1)},\ldots,X_{(n)})\in B)$, 这正是均匀性与独立性.
\end{proof}

第 $k$ 小的密度为
\[f_{(k)}(x)=\frac{n!}{(k-1)!(n-k)!}F(x)^{k-1}(1-F(x))^{n-k}f(x).\]
理由是选出一个样本落在 $[x,x+\dd x]$, 选出 $k-1$ 个在其左边, 剩余的在右边; 有 $n\binom{n-1}{k-1}$ 种编号选择. 也可把联合排序密度的其余坐标积分掉, 左侧有序积分为 $F(x)^{k-1}/(k-1)!$, 右侧为 $(1-F(x))^{n-k}/(n-k)!$, 得到同一公式.
\begin{example}
独立抽取七个 $(0,1)$ 均匀样本. 已知第一个样本在七个中位列第五小, 求它的条件密度、平均值及众数.
\end{example}
\begin{solution}
名次排列与排序值独立, 所求分布是 $X_{(5)}$ 的分布, 密度为 $105x^4(1-x)^2$, $0<x<1$. 一般整数参数的贝塔积分
\[\int_0^1x^{a-1}(1-x)^{b-1}\dd x=\frac{(a-1)!(b-1)!}{(a+b-1)!}\]
由分部积分递推得到: 当 $b>1$ 时, 对 $x^{a-1}$ 积分, 将 $(1-x)^{b-1}$ 求导, 得原积分等于 $(b-1)/a$ 倍参数为 $(a+1,b-1)$ 的积分, 最后归结为 $\int_0^1x^{a+b-2}\dd x=1/(a+b-1)$. 因此平均值是相邻贝塔积分之比 $a/(a+b)$, 本题 $a=5,b=3$, 平均值为 $5/8$. 密度对数的导数是 $4/x-2/(1-x)$, 在 $x=2/3$ 为0, 左正右负, 所以众数为 $2/3$. 一般 $a,b>1$ 时众数为 $(a-1)/(a+b-2)$.
\end{solution}

\originalsection{协方差与和的方差}
\begin{proposition}[独立变量的乘积]
若 $X,Y$ 独立且 $g(X),h(Y)$ 可积, 则乘积也可积, 且
$\E[g(X)h(Y)]=\E g(X)\E h(Y)$.
\end{proposition}
\begin{proof}
离散情形把联合质量分解为 $p_X(x)p_Y(y)$, 连续情形把联合密度分解为 $f_X(x)f_Y(y)$. 先对绝对值求和或积分得到 $\E|g(X)h(Y)|=\E|g(X)|\E|h(Y)|<\infty$, 再交换求和或积分并分离两个因子. 一般独立分布也有相同的乘积分解.
\end{proof}
\begin{definition}[协方差]
若 $X,Y$ 有有限二阶矩, 定义
\[\Cov(X,Y)=\E[(X-\E X)(Y-\E Y)].\]
有限二阶矩使乘积可积, 因为 $2|ab|\le a^2+b^2$. 展开得
$\Cov(X,Y)=\E(XY)-\E X\E Y$.
\end{definition}
协方差为正意味着两个相对各自平均值的偏差, 同号的贡献超过异号的贡献. 它不衡量一切形式的依赖, 只衡量这种乘积平均.
\begin{proposition}[双线性与方差公式]
协方差对交换变量对称, $\Cov(X,X)=\Var X$, 且
\[
\Cov\left(\sum_i a_iX_i,\sum_j b_jY_j\right)
=\sum_{i,j}a_ib_j\Cov(X_i,Y_j).
\]
因此
\[\Var\left(\sum_iX_i\right)=\sum_i\Var X_i+2\sum_{i<j}\Cov(X_i,X_j).\]
\end{proposition}
\begin{proof}
减去各自期望后, 第一个和变成 $\sum_i a_i(X_i-\E X_i)$, 第二个和相同处理. 相乘并用有限线性性便得双线性式. 令两组变量相同且系数全为1, 对角项是方差, 每一非对角项出现两次. 独立时乘积命题说明各交叉协方差为0; 两两不相关其实已经足够使方差相加.
\end{proof}
\begin{example}
独立的 $X_1,\ldots,X_5$ 方差均为 $v>0$. 令 $A=X_1+X_2+X_3+X_4$, $B=X_3+X_4+X_5$. 求协方差与相关系数.
\end{example}
\begin{solution}
展开协方差后, 仅重复出现的 $X_3,X_4$ 贡献 $v$, 所以 $\Cov(A,B)=2v$. 又 $\Var A=4v$, $\Var B=3v$, 相关系数为 $2/\sqrt{12}=1/\sqrt3$. 一般独立同方差变量的两个子集和, 协方差等于交集大小乘共同方差.
\end{solution}
\begin{example}
继续包裹随机分配问题, 设 $m\ge2$. 求送对数 $N$ 的方差.
\end{example}
\begin{solution}
$\Var I_i=(1/m)(1-1/m)$. 对 $i\ne j$, 同时送对这两个包裹的排列有 $(m-2)!$ 个, 故 $\E(I_iI_j)=1/[m(m-1)]$, 从而
\[\Cov(I_i,I_j)=\frac1{m(m-1)}-\frac1{m^2}=\frac1{m^2(m-1)}.\]
代入和的方差公式,
\[\Var N=m\frac1m\left(1-\frac1m\right)+m(m-1)\frac1{m^2(m-1)}=1.\]
这些指标并不独立, 忽略协方差会少算 $1/m$. $m=1$ 时 $N$ 恒等于1, 方差为0, 不属于上述计算范围.
\end{solution}

\originalsection{相关系数及其局限}
\begin{definition}
当 $0<\Var X,\Var Y<\infty$ 时, 定义
$\rho(X,Y)=\Cov(X,Y)/\sqrt{\Var X\Var Y}$. 常量变量的方差为0, 其相关系数不定义.
\end{definition}
\begin{theorem}
$-1\le\rho(X,Y)\le1$. 等号 $|\rho|=1$ 当且仅当存在 $a\ne0,b\in\R$, 使 $Y=aX+b$ 几乎必然成立. 此时 $\rho$ 等于 $a$ 的符号.
\end{theorem}
\begin{proof}
令 $U=(X-\E X)/\sqrt{\Var X}$, $V=(Y-\E Y)/\sqrt{\Var Y}$. 二者均值为0、二阶矩为1, 且 $\E UV=\rho$. 于是 $\E(U-V)^2=2-2\rho\ge0$, $\E(U+V)^2=2+2\rho\ge0$, 得上下界. 若 $\rho=1$, 非负变量 $(U-V)^2$ 平均为0, 故 $U=V$ 几乎必然; 若 $\rho=-1$ 则 $U=-V$. 两种情况都转成所述仿射关系. 反向代入协方差定义即可验证.
\end{proof}
对 $a,c\ne0$, $\rho(aX+b,cY+d)=\operatorname{sgn}(ac)\rho(X,Y)$. 正比例改变单位不会改变相关系数, 负比例会翻转方向. 独立且方差有限、非零的变量一定不相关, 逆命题不成立.
\begin{example}
令 $X$ 在 $(-1,1)$ 均匀分布, $Y=X^2$. 证明它们不相关但不独立.
\end{example}
\begin{solution}
对称性给 $\E X=\E X^3=0$, 所以 $\Cov(X,Y)=0$. 两变量方差均为正. 但事件 $A=\{|X|<1/2\}$ 与 $B=\{Y<1/4\}$ 相同, 概率为 $1/2$. 因而 $\Prob(A\cap B)=1/2\ne1/4=\Prob(A)\Prob(B)$, 它们不独立. 相关为0只排除了线性共同变化, 没有排除平方关系.
\end{solution}

\originalsection{无限平均值与有限决策}
有限线性性中的可积条件不能随意删除. $\infty-\infty$ 没有数值意义, 无限序列中每一步改善也不保证极限对象改善. 例如起初持有编号1的凭证, 第 $n$ 步交出编号 $n$ 的凭证并收到编号 $2n,2n+1$ 的凭证. 每一步净增一个, 完成 $n$ 步时持有 $n+1$ 个; 但固定编号的凭证最终都会交出, 所以逐编号极限为空集. 计数与这一极限不交换, 不能把有限阶段的净收益直接延伸到无穷阶段.

同样, 对每个有限 $n$, 先等待 $n$ 天再以 $1/n$ 的概率遭遇无限等待的策略, 都有无限期望损失. 把 $1/n$ 降为 $1/(n+1)$ 的好处乘以无穷, 再据此无限推迟, 并不是两个有限期望之间的比较. 若损失或效用无界, 应先说明采用哪一种有限截断或优化准则.

无限总量的搬移也会产生类似错觉. 若第 $n$ 堆有 $r^n$ 单位, $r>1$, 把比例 $q$ 移向左边, 剩余移向右边, 则内部第 $n$ 堆的新量为 $q r^{n+1}+(1-q)r^{n-1}$. 例如 $r=4,q=1/2$, 每堆变为原来的 $17/8$ 倍. 这没有从有限总量制造财富: 初始总量已经无穷, 边界流入流出和截断总量必须另算.
\begin{exercise}
令 $K$ 满足 $\Prob(K=k)=2^{-(k+1)}$, $k=0,1,\ldots$. 两个信封的金额为 $A=3^K$ 与 $2A$, 随机等可能打开一个, 看到金额 $V$. 计算看到 $V=1$ 和看到其他可能金额时换到另一个信封的条件平均收益, 并解释为何不能推出无条件平均收益为正.
\end{exercise}
\begin{solution}
幂 $3^k$ 与 $2\cdot3^j$ 不会相等. 所以看到 $V=3^k$ 就知道自己拿的是小信封, 换后收益为 $3^k$; 看到 $V=2\cdot3^k$ 就知道是大信封, 收益为 $-3^k$. 特别地 $V=1$ 时收益为1. 无条件收益的正、负部分期望分别为
$\frac12\sum_{k\ge0}3^k2^{-(k+1)}=\infty$, 因此收益的期望未定义, 不能通过把正、负无穷相减比较策略. 已知金额后的判断必须使用实际条件分布, 不能不加检查地说另一个金额以各半概率是当前金额的两倍或一半.
\end{solution}
